% !TeX root = ../../main.tex

Retry-Mechanismen stellen eine verbreitete Strategie dar, um temporäre Fehler automatisiert zu behandeln.
Besonders geeignet sind sie für Fehler, deren Ursache sich innerhalb eines kurzen Zeitraums selbstständig behebt.
Dazu zählen Verbindungsabbrüche, Überlastungen oder Konflikte durch parallele Transaktionen.
In solchen Fällen kann eine erneute Ausführung dazu führen, dass eine Transaktion erfolgreich abschließt ohne
einen manuellen Eingriff.

Da Retry eine Transaktion einfach nur wiederholt, wird dadurch das System mehr belastet.
Deswegen muss abgewägt werden, ob ein schneller Abbruch der Transaktion nicht sinnvoller ist, als häufiges Wiederholen.
Des Weiteren kann vor allem aggressives Wiederholen, also häufiges Wiederholen mit geringen Intervallen, dafür
sorgen, dass ein System an die Belastungsgrenze stößt.
Das ist vor allem ein Problem für Systeme, die von Anfang an am Ende ihrer Kapazitäten laufen.
Um unnötiges Wiederholen der Transaktionen zu vermeiden, sollten die Fehlercodes ausgelesen werden, um
sicherzustellen, dass es sich um ein transient failure handelt und nicht um einen anderen Fehler.
Darüber hinaus sollte ein Retry-Mechanismus nur für idempotente oder entsprechend abgesicherte Transaktionen
eingesetzt werden~\cite{MicrosoftRetryPattern}.
Außerdem müssen die Anzahl der Retries und der Delay optimal gewählt werden, denn sonst besteht die Möglichkeit,
dass ein Retry zum einen keinen Effekt hat oder zu viele Ressourcen verbraucht~\cite{MicrosoftTransientFaults}.

Wie wirksam ein Retry ist, hängt stark von der Art des aufgetretenen Fehlers ab.
Bei Fehlern mit dauerhafter Ursache kann eine Wiederholung keine Verbesserung erzielen.
Ein solcher Fehler kann entstehen, wenn die Transaktion inhaltliche Fehler enthält.
In diesem Fall würde ein erneuter Versuch das gleiche Ergebnis liefern.
Die Fehlerbehebung muss manuell erfolgen.
Zusätzlich hängt der Erfolg eines Retries von der gewählten Strategie ab.
Eine schlecht optimierte Konfiguration kann die Situation verschärfen, anstatt sie zu beheben.

Neben den Grenzen der Fehlerbehebung können Retry-Mechanismen selbst neue Herausforderungen erzeugen.
Jeder wiederholte Versuch einer Transaktion verursacht zusätzliche Last für die Datenbank.
Besonders bei hoher Anzahl paralleler Anfragen kann dies zu verstärkter Konkurrenz um Ressourcen führen.
Dadurch können sich Antwortzeiten verlängern und die Leistung des Systems verschlechtern.
Insbesondere bei vielen gleichzeitig ausgeführten Retries kann dadurch ein Kreislauf entstehen, bei dem zusätzliche
Datenbanklast wiederum neue Fehler verursacht.

Darüber hinaus besteht keine Garantie, dass eine wiederholte Transaktion erfolgreich abgeschlossen wird.
Wenn Retries ohne passende Optimierung eingesetzt werden, können zahlreiche Wiederholungsversuche entstehen.
Diese verbrauchen Ressourcen, ohne dass es einen Nutzen hat.
Daher ist eine geeignete Retry-Strategie notwendig, die sowohl die Fehlerart als auch die Auswirkungen auf die
Datenbank berücksichtigt.